\begin{frame}{Satisfiability}
\begin{center}
\vspace{1.5cm}
{\Huge\textbf{Satisfiability}}
\end{center}
\end{frame}

% --------------------------------------------------
% Question 1
% --------------------------------------------------

\begin{frame}{Question 1}
A software system has the following access rules:

\begin{itemize}
\item A user can access the system if the password is valid.
\item A user cannot access the system if the password is invalid.
\end{itemize}

Let:

$$
P = \text{Password is valid}
$$

$$
Q = \text{User can access the system}
$$

Determine whether the following set of conditions is satisfiable:

$$
P \rightarrow Q
$$

$$
\neg P \rightarrow \neg Q
$$

\end{frame}

\begin{frame}{Question 1 -- Answer}
Consider the possible values of $P$.

If:

$$
P=T
$$

then the first condition requires:

$$
Q=T
$$

If:

$$
P=F
$$

then the second condition requires:

$$
Q=F
$$

\end{frame}

\begin{frame}{Question 1 -- Answer}
A satisfying assignment is therefore:

$$
P=T,\qquad Q=T
$$

Both conditions are true:

$$
T\rightarrow T = T
$$

$$
F\rightarrow F = T
$$

Therefore:

$$
\boxed{\text{The set of conditions is satisfiable.}}
$$

\end{frame}

% --------------------------------------------------
% Question 2
% --------------------------------------------------

\begin{frame}{Question 2}
A database system has the following requirements:

\begin{itemize}
\item A transaction must be approved before execution.
\item The transaction is executed.
\item The transaction is not approved.
\end{itemize}

Let:

$$
P = \text{Transaction is approved}
$$

$$
Q = \text{Transaction is executed}
$$

Determine whether the following conditions are satisfiable:
$$
P\rightarrow Q
$$
$$
Q
$$
$$
\neg P
$$

\end{frame}

\begin{frame}{Question 2 -- Answer}
The conditions require:

$$
Q=T
$$

and:

$$
P=F
$$

Substituting these values into the implication gives:

$$
F\rightarrow T = T
$$

\end{frame}

\begin{frame}{Question 2 -- Answer}
Therefore, the assignment

$$
P=F,\qquad Q=T
$$

satisfies all three conditions.

Thus:

$$
\boxed{\text{The conditions are satisfiable.}}
$$

The implication does not require $P$ to be true whenever $Q$ is true.
\end{frame}

% --------------------------------------------------
% Question 3
% --------------------------------------------------

\begin{frame}{Question 3}
An authentication system has the following rules:

\begin{itemize}
\item If a user is authenticated, access is granted.
\item Access is granted.
\item The user is not authenticated.
\end{itemize}

Let:

$$
P = \text{User is authenticated}
$$

$$
Q = \text{Access is granted}
$$

Determine whether the following set of statements is satisfiable:
$$
P\rightarrow Q
$$
$$
Q
$$
$$
\neg P
$$

\end{frame}

\begin{frame}{Question 3 -- Answer}
The conditions require:

$$
Q=T
$$

and:

$$
P=F
$$

The implication becomes:

$$
F\rightarrow T=T
$$

\end{frame}

\begin{frame}{Question 3 -- Answer}
Therefore:

$$
P=F,\qquad Q=T
$$

makes all the statements true.

Hence:

$$
\boxed{\text{The statements are satisfiable.}}
$$

This example shows that an implication alone does not allow us to infer
that its antecedent is true.
\end{frame}

% --------------------------------------------------
% Question 4
% --------------------------------------------------

\begin{frame}{Question 4}
A cloud service has the following requirements:

\begin{itemize}
\item A server is operational if and only if it is connected to the
network.
\item The server is operational.
\item The server is not connected to the network.
\end{itemize}

Let:

$$
P = \text{Server is operational}
$$

$$
Q = \text{Server is connected to the network}
$$

Determine whether the requirements are satisfiable.
\end{frame}

\begin{frame}{Question 4 -- Answer}
The first requirement is:

$$
P\leftrightarrow Q
$$

The other requirements state:

$$
P=T
$$

and:

$$
Q=F
$$

\end{frame}

\begin{frame}{Question 4 -- Answer}
Substituting these values into the biconditional gives:

$$
T\leftrightarrow F=F
$$

Therefore, the first requirement is violated.

There is no assignment that can satisfy all three requirements.
\end{frame}

\begin{frame}{Question 4 -- Answer}
Hence:

$$
\boxed{\text{The requirements are unsatisfiable.}}
$$

The contradiction occurs because the biconditional requires $P$ and $Q$
to have the same truth value.
\end{frame}

% --------------------------------------------------
% Question 5
% --------------------------------------------------

\begin{frame}{Question 5}
A machine-learning deployment system uses the following policy:

\begin{itemize}
\item A model is deployed only if it passes testing.
\item A model is deployed only if its accuracy is at least $90\%$.
\item The model is deployed.
\item The model does not pass testing.
\end{itemize}

\end{frame}

\begin{frame}{Question 5}

Let:
$$
P = \text{Model passes testing}
$$
$$
Q = \text{Model accuracy is at least }90\%
$$
$$
R = \text{Model is deployed}
$$

Determine whether the following set of conditions is satisfiable:

$$
P\rightarrow R
$$
$$
Q\rightarrow R
$$
$$
R
$$

$$
\neg P
$$

\end{frame}

\begin{frame}{Question 5 -- Answer}
The given conditions require:

$$
R=T
$$

and:

$$
P=F
$$

Consider the first implication:

$$
P\rightarrow R
$$

Substituting the values gives:

$$
F\rightarrow T=T
$$

\end{frame}

\begin{frame}{Question 5 -- Answer}
The second implication is:

$$
Q\rightarrow R
$$

Since $R=T$, this implication is true regardless of the value of $Q$.

For example:

$$
Q=T
$$

gives:

$$
T\rightarrow T=T
$$

\end{frame}

\begin{frame}{Question 5 -- Answer}
A satisfying assignment is:

$$
P=F,\qquad Q=T,\qquad R=T
$$

All four conditions are satisfied.

Therefore:

$$
\boxed{\text{The set of conditions is satisfiable.}}
$$

\end{frame}
